\documentclass[10pt]{article}

%----------------------------------------------------------
% Geometry and layout
%----------------------------------------------------------
\usepackage[a4paper,margin=2.5cm]{geometry}
\usepackage{fancyhdr}
\usepackage{graphicx}
\usepackage{booktabs}
\usepackage{amsmath,amssymb}
\usepackage{hyperref}
\usepackage[labelfont=bf]{caption}

%----------------------------------------------------------
% Page style: Supplementary Material header + bottom-center page numbers
%----------------------------------------------------------
\pagestyle{fancy}
\fancyhf{}
\fancyhead[C]{\textsc{Supplementary Material}}
\renewcommand{\headrulewidth}{0.4pt}
\fancyfoot[C]{\thepage}
\renewcommand{\footrulewidth}{0pt}

% Plain style for the title page
\fancypagestyle{plain}{
  \fancyhf{}
  \fancyhead[C]{\textsc{Supplementary Material}}
  \renewcommand{\headrulewidth}{0.4pt}
  \fancyfoot[C]{\thepage}
}

%----------------------------------------------------------
% Caption setup: prefix all figure/table labels with "Supplementary"
%----------------------------------------------------------
\makeatletter
\renewcommand{\fnum@figure}{Supplementary Figure~\thefigure}
\renewcommand{\fnum@table}{Supplementary Table~\thetable}
\makeatother

%----------------------------------------------------------
% Metadata
%----------------------------------------------------------
\title{\textbf{Supplementary Material for:}\\[0.5em]
\Large Forty Years of Interatomic Potentials, One Error Geometry:
Hyper-Ribbon Universality and Why Pooled Benchmarks Mislead}
\author{Alex Welcing}
\date{}

%----------------------------------------------------------
\begin{document}
\maketitle
\thispagestyle{plain}

\noindent
This document contains supplementary figures referenced in the main article. 
The main manuscript presents six figures (the journal limit); the material below 
provides additional analyses that support the central claims without overloading 
the primary narrative.

\vspace{1em}

%----------------------------------------------------------
% Supplementary Figure 1
%----------------------------------------------------------
\clearpage
\begin{figure}[p]
    \centering
    \includegraphics[width=0.85\textwidth]{figures/year_stratified_dim.png}
    \caption{Effective dimensionality (Participation Ratio) of prediction errors
    stratified by potential publication year (1980--2025). The hyper-ribbon
    structure remains stable across all decades, with PR bounded well below the
    ambient dimension 3 (median $\approx 1.1$) regardless of model complexity or
    publication era.}
    \label{fig:supp1}
\end{figure}

%----------------------------------------------------------
% Supplementary Figure 2
%----------------------------------------------------------
\clearpage
\begin{figure}[p]
    \centering
    \includegraphics[width=0.85\textwidth]{figures/year_stratified_r2.png}
    \caption{Log-eigenvalue linearity ($R^2$) stratified by publication year. 
    The Vandermonde universality class signature ($R^2 > 0.80$) persists across 
    all decades of interatomic potential development.}
    \label{fig:supp2}
\end{figure}

%----------------------------------------------------------
% Supplementary Figure 3
%----------------------------------------------------------
\clearpage
\begin{figure}[p]
    \centering
    \includegraphics[width=0.85\textwidth]{figures/observables_5d_pr.png}
    \caption{Hyper-ribbon dimensionality distribution comparing the standard 3D 
    elastic constant space ($C_{11}$, $C_{12}$, $C_{44}$) to an expanded 5D space 
    incorporating lattice constant ($a$) and cohesive energy ($E_{\text{coh}}$). 
    The median effective dimensionality remains highly compressed despite a 66\% 
    increase in the number of target observables, confirming that the low-dimensional 
    error structure is fundamental to the observable space rather than an artifact 
    of dimensionality restriction.}
    \label{fig:supp3}
\end{figure}

\end{document}
